\section{Notation, parameters and initial conditions}
Electrical powers are in kW, thermal powers in kW-th, stored energies in kWh or kWh-th, temperatures in degrees Celsius, and computational work in aggregate CPU-hours. Thermal heat capacities and conductances below use kJ/K and kW/K. The implementation uses W and J for several cooling quantities; the equations here convert these consistently to kW and kJ. A rate of one denotes the entire modelled computing capacity.

\begin{longtable}{@{}p{0.24\textwidth}p{0.69\textwidth}@{}}
\caption{Principal notation.}\label{tab:symbols}\\\toprule
Symbol & Meaning\\\midrule\endfirsthead
\toprule Symbol & Meaning\\\midrule\endhead
$t$, $\Delta t$, $\Delta t_s$ & Interval start; interval length 0.25 h or 900 s. Boundary $t+1$ follows interval $t$.\\
$d$, $\mathcal C_d$, $\mathcal H_d$ & Local date, committed core intervals, core plus twelve look-ahead intervals.\\
$c$, $r_c$, $\ell_c$, $a_c$ & Cohort, arrival, final eligible execution-interval start, original CPU-hour requirement.\\
$A_{c,d}$, $R_{c,H}$ & Work available to process at horizon opening (or arrival within it), and remaining work at boundary $H$.\\
$u^{\mathrm{inflex}}_t$, $u^{\mathrm{flex,arr}}_t$, $u_{c,t}$, $u_t$ & Inflexible and flexible arrival rates, scheduled cohort rate, total executed utilisation.\\
$\alpha_{t,k}$, $D_k$ & Fraction of flexible arrival assigned to class $k$, and its maximum shift in interval starts.\\
$P^{\mathrm{IT}}_t$, $P^{\mathrm{grid}}_t$, $P^{\mathrm{aux}}$ & IT power, site import, fixed auxiliary power.\\
$P^{\mathrm{U,ch}}_t$, $P^{\mathrm{U,dis}}_t$, $E^{\mathrm U}_t$, $z_t^{\mathrm U}$ & UPS charge/discharge power, stored energy and charge-mode binary.\\
$P^{\mathrm C}_t$, $P^{\mathrm{T,ch}}_t$ & Chiller electricity for direct CRAC cooling and TES charging.\\
$Q^{\mathrm{T,ch}}_t$, $Q^{\mathrm{T,dis}}_t$, $E^{\mathrm T}_t$, $z_t^{\mathrm T}$ & TES thermal charge/discharge, stored cooling energy and charge-mode binary.\\
$Q_t^{\mathrm{cool}}$ & Gross cooling supply from direct chiller and TES paths.\\
$T^{\mathrm{IT}}$, $T^{\mathrm R}$, $T^{\mathrm{CA}}$, $T^{\mathrm{HA}}$, $T^{\mathrm{in}}$ & IT, rack, cold-aisle and hot-aisle state temperatures; interval inlet temperature.\\
$C_j$, $G_{\mathrm{conv}}$, $G_{\mathrm{cold}}$, $\dot m$, $c_p$, $\kappa$ & Thermal capacities, IT--rack and cold-aisle conductances, air flow, specific heat and effective-airflow factor.\\
$\pi_t$, $C^{\mathrm{ref}}$, $C^{\op}$ & Signed reference price; annual comparator and coordinated costs.\\
$t_0$, $\Delta P$, $\tau$, $H$ & Event start, signed import deviation, event duration and recovery boundary.\\
Superscripts $\mathrm{ref}$, $\op$ & Workload-at-arrival comparator and undisturbed coordinated operation.\\\bottomrule
\end{longtable}

\begin{table}[htbp]
\centering\small
\caption{Thermal constants from the implemented parameter definitions (rounded here to six decimal places).}\label{tab:thermal}
\begin{tabular}{@{}lrl@{}}\toprule
Parameter & Value & Unit \\\midrule
$C_{\mathrm{IT}}$ & 17880.000000 & kJ/K \\
$C_{\mathrm R}$ & 18019.516902 & kJ/K \\
$C_{\mathrm{CA}}$ & 2332.644000 & kJ/K \\
$C_{\mathrm{HA}}$ & 1166.322000 & kJ/K \\
$G_{\mathrm{conv}}$ & 108.991826 & kW/K \\
$G_{\mathrm{cold}}$ & 4.483672 & kW/K \\
$c_p$ & 1.005450 & kJ/(kg K) \\
$\dot m$ & 100.000000 & kg/s \\
$\kappa$ & 0.766300 & -- \\
$T^{\mathrm{out}}$ & 22.000000 & $^{\circ}$C \\
\bottomrule\end{tabular}
\end{table}


The principal electrical and storage settings are: IT idle/rated power 166.7/1,000 kW; UPS energy 600 kWh, bounds 300--600 kWh, initial energy 600 kWh; UPS charge limit 270 kW, effective discharge limit 1,000 kW, efficiencies 0.82/0.92; TES bounds 0--1,000 kWh-th, initial energy 500 kWh-th, charge/discharge limits 300 kW-th, efficiencies 0.9/0.9; chiller input limit 400 kW and COP 5. The UPS's nominal discharge rating is 2,700 kW, but the model restricts useful discharge to IT demand and hence to at most 1,000 kW. The time-varying IT restriction can be tighter still.

Initial temperatures $(T^{\mathrm{IT}},T^{\mathrm R},T^{\mathrm{CA}},T^{\mathrm{HA}})$ are $(28.5,26,20,27)^{\circ}$C. Bounds at every state boundary are 18--60, 18--40, 18--23 and 18--40$^{\circ}$C, respectively; inlet temperature is 18--30$^{\circ}$C. Both annual cases use the same thermal bounds. A historical configuration field named \path{baseline_cold_aisle_setpoint_c} is not an active fixed-setpoint constraint in the retained reference model.

Auxiliary power is 53.095 kW. It was calibrated as $0.07\times758.5$ kW from an earlier daily IT-plus-cooling benchmark and is held constant in all retained cases. It represents approximately 7.4\% of annual reference IT-plus-cooling mean demand, rather than an enforced 7\% annual share. The structural grid-import limit is $1000+270+400+53.095=1723.095$ kW. No additional site-specific connection cap, ramp-rate constraint or reserve above the proportional UPS lower bound is active in these results.

\section{Workload and electricity-price inputs}
\subsection{Adopted workload arrays and their interpretation}
Published studies describe differing utilisation levels, workload classes and temporal flexibility across facilities~\cite{Cao2022ApEn,6779441,Misaghian2023TIA,Liu2012PER}. They motivate separating immediate work from work that may be rescheduled, but do not identify a universal flexible/inflexible ratio. The inputs used here are stylised profiles selected with reference to that literature. They are not claimed to be a measured site trace, a statistically representative sample of facilities or a deterministic reconstruction of a particular published trace.

The numerical inputs are \path{data/load_profiles.csv} and \path{data/shiftability_profile.csv}, whose active slots 1--96 are reproduced jointly in Table~\ref{tab:workload-inputs}. The load file specifies flexible and inflexible arrival utilisation at quarter-hour resolution. It also contains legacy rows 97--108; annual timeline construction does not use those rows. Utilisation may vary within an hour: for example, flexible utilisation at 00:00, 00:15, 00:30 and 00:45 is 0.40, 0.38, 0.36 and 0.34. Thus the hourly starting value is 0.40, whereas the hourly mean is 0.37. The model uses the four individual values, without replacing them by an hourly mean or a constant hourly starting value.

Tranche shares have a different construction. Table~\ref{tab:hourly-tranches} specifies the intended 24 hourly distributions. Each hourly row is repeated four consecutive times to produce the 96-row input: the 00:00--01:00 distribution applies at 00:00, 00:15, 00:30 and 00:45, followed by four copies of the 01:00--02:00 distribution, and so on. Each distribution sums to one and allocates the flexible work arriving in that interval. The 24-row sequence is not repeated four times across the day. The supplied quarter-hour CSV and Table~\ref{tab:workload-inputs} reproduce this expansion exactly.

\begin{table}[htbp]
\centering\small
\caption{Intended hourly allocation of flexible arrivals (percent). Each row is held constant over its four quarter-hour intervals.}\label{tab:hourly-tranches}
\begin{tabular}{@{}lrrrr@{}}\toprule
Local hour & $\alpha_1$ & $\alpha_2$ & $\alpha_3$ & $\alpha_4$ \\\midrule
00:00--01:00 & 25 & 25 & 20 & 30 \\
01:00--02:00 & 25 & 25 & 15 & 35 \\
02:00--03:00 & 35 & 17 & 18 & 30 \\
03:00--04:00 & 25 & 15 & 20 & 40 \\
04:00--05:00 & 25 & 20 & 15 & 40 \\
05:00--06:00 & 25 & 15 & 15 & 45 \\
06:00--07:00 & 23 & 28 & 15 & 34 \\
07:00--08:00 & 20 & 20 & 15 & 45 \\
08:00--09:00 & 25 & 15 & 13 & 47 \\
09:00--10:00 & 25 & 17 & 16 & 42 \\
10:00--11:00 & 35 & 22 & 13 & 30 \\
11:00--12:00 & 32 & 20 & 15 & 33 \\
12:00--13:00 & 45 & 25 & 15 & 15 \\
13:00--14:00 & 50 & 20 & 15 & 15 \\
14:00--15:00 & 40 & 20 & 25 & 15 \\
15:00--16:00 & 45 & 25 & 15 & 15 \\
16:00--17:00 & 50 & 15 & 20 & 15 \\
17:00--18:00 & 50 & 20 & 15 & 15 \\
18:00--19:00 & 10 & 20 & 30 & 40 \\
19:00--20:00 & 15 & 20 & 25 & 40 \\
20:00--21:00 & 18 & 12 & 25 & 45 \\
21:00--22:00 & 22 & 18 & 25 & 35 \\
22:00--23:00 & 20 & 16 & 24 & 40 \\
23:00--24:00 & 15 & 20 & 20 & 45 \\
\bottomrule\end{tabular}
\end{table}


For an arrival in interval $t$, class $k$ receives work
\begin{equation}
a_{t,k}=u_t^{\mathrm{flex,arr}}\alpha_{t,k}\Delta t,
\quad \ell_{t,k}=t+D_k\Delta t,
\quad (D_1,D_2,D_3,D_4)=(2,4,8,12).
\end{equation}
The final execution interval starts at $\ell_{t,k}$ and closes at $\ell_{t,k}+\Delta t$. For example, a three-hour class arriving at 12:00 may execute in the interval starting at 15:00 and ending at 15:15. Thus the class names denote shifts between interval starts, not a guarantee of completion within exactly three continuous hours of an interval's opening. Cohort creation omits quantities no larger than the $10^{-7}$ CPU-hour numerical tolerance.

For workload multiplier $m$, let $v_t=u_t^{\mathrm{inflex}}+u_t^{\mathrm{flex,arr}}$ in the original input. The sensitivity uses
\begin{equation}
u_t^{\mathrm{flex,arr}}(m)=\min\{m u_t^{\mathrm{flex,arr}}(1),v_t\},
\qquad u_t^{\mathrm{inflex}}(m)=v_t-u_t^{\mathrm{flex,arr}}(m).
\end{equation}
Tranche fractions remain fixed. The clipping rule preserves total arriving CPU-hours and avoids creating more flexible work than the total workload. A multiplier is therefore not always an exactly proportional increase in the realised flexible share.

\begin{longtable}{@{}lrrrrrr@{}}
\caption{Actual quarter-hour workload inputs. Utilisations are fractions of rated CPU capacity; tranche shares are fractions of flexible work.}\label{tab:workload-inputs}\\\toprule
Local start & Inflexible & Flexible & $\alpha_1$ & $\alpha_2$ & $\alpha_3$ & $\alpha_4$ \\\midrule\endfirsthead
\toprule
Local start & Inflexible & Flexible & $\alpha_1$ & $\alpha_2$ & $\alpha_3$ & $\alpha_4$ \\\midrule\endhead
00:00 & 0.28 & 0.40 & 0.25 & 0.25 & 0.20 & 0.30 \\
00:15 & 0.27 & 0.38 & 0.25 & 0.25 & 0.20 & 0.30 \\
00:30 & 0.26 & 0.36 & 0.25 & 0.25 & 0.20 & 0.30 \\
00:45 & 0.25 & 0.34 & 0.25 & 0.25 & 0.20 & 0.30 \\
01:00 & 0.25 & 0.31 & 0.25 & 0.25 & 0.15 & 0.35 \\
01:15 & 0.23 & 0.32 & 0.25 & 0.25 & 0.15 & 0.35 \\
01:30 & 0.21 & 0.32 & 0.25 & 0.25 & 0.15 & 0.35 \\
01:45 & 0.19 & 0.32 & 0.25 & 0.25 & 0.15 & 0.35 \\
02:00 & 0.17 & 0.33 & 0.35 & 0.17 & 0.18 & 0.30 \\
02:15 & 0.16 & 0.30 & 0.35 & 0.17 & 0.18 & 0.30 \\
02:30 & 0.16 & 0.28 & 0.35 & 0.17 & 0.18 & 0.30 \\
02:45 & 0.16 & 0.26 & 0.35 & 0.17 & 0.18 & 0.30 \\
03:00 & 0.16 & 0.23 & 0.25 & 0.15 & 0.20 & 0.40 \\
03:15 & 0.14 & 0.24 & 0.25 & 0.15 & 0.20 & 0.40 \\
03:30 & 0.12 & 0.25 & 0.25 & 0.15 & 0.20 & 0.40 \\
03:45 & 0.10 & 0.26 & 0.25 & 0.15 & 0.20 & 0.40 \\
04:00 & 0.08 & 0.27 & 0.25 & 0.20 & 0.15 & 0.40 \\
04:15 & 0.07 & 0.27 & 0.25 & 0.20 & 0.15 & 0.40 \\
04:30 & 0.06 & 0.27 & 0.25 & 0.20 & 0.15 & 0.40 \\
04:45 & 0.06 & 0.27 & 0.25 & 0.20 & 0.15 & 0.40 \\
05:00 & 0.06 & 0.27 & 0.25 & 0.15 & 0.15 & 0.45 \\
05:15 & 0.07 & 0.24 & 0.25 & 0.15 & 0.15 & 0.45 \\
05:30 & 0.09 & 0.22 & 0.25 & 0.15 & 0.15 & 0.45 \\
05:45 & 0.10 & 0.20 & 0.25 & 0.15 & 0.15 & 0.45 \\
06:00 & 0.12 & 0.18 & 0.23 & 0.28 & 0.15 & 0.34 \\
06:15 & 0.14 & 0.19 & 0.23 & 0.28 & 0.15 & 0.34 \\
06:30 & 0.16 & 0.21 & 0.23 & 0.28 & 0.15 & 0.34 \\
06:45 & 0.18 & 0.22 & 0.23 & 0.28 & 0.15 & 0.34 \\
07:00 & 0.20 & 0.24 & 0.20 & 0.20 & 0.15 & 0.45 \\
07:15 & 0.21 & 0.24 & 0.20 & 0.20 & 0.15 & 0.45 \\
07:30 & 0.22 & 0.24 & 0.20 & 0.20 & 0.15 & 0.45 \\
07:45 & 0.23 & 0.24 & 0.20 & 0.20 & 0.15 & 0.45 \\
08:00 & 0.24 & 0.24 & 0.25 & 0.15 & 0.13 & 0.47 \\
08:15 & 0.26 & 0.25 & 0.25 & 0.15 & 0.13 & 0.47 \\
08:30 & 0.29 & 0.26 & 0.25 & 0.15 & 0.13 & 0.47 \\
08:45 & 0.31 & 0.27 & 0.25 & 0.15 & 0.13 & 0.47 \\
09:00 & 0.34 & 0.28 & 0.25 & 0.17 & 0.16 & 0.42 \\
09:15 & 0.35 & 0.25 & 0.25 & 0.17 & 0.16 & 0.42 \\
09:30 & 0.36 & 0.23 & 0.25 & 0.17 & 0.16 & 0.42 \\
09:45 & 0.37 & 0.21 & 0.25 & 0.17 & 0.16 & 0.42 \\
10:00 & 0.37 & 0.19 & 0.35 & 0.22 & 0.13 & 0.30 \\
10:15 & 0.38 & 0.18 & 0.35 & 0.22 & 0.13 & 0.30 \\
10:30 & 0.39 & 0.17 & 0.35 & 0.22 & 0.13 & 0.30 \\
10:45 & 0.40 & 0.16 & 0.35 & 0.22 & 0.13 & 0.30 \\
11:00 & 0.42 & 0.16 & 0.32 & 0.20 & 0.15 & 0.33 \\
11:15 & 0.41 & 0.17 & 0.32 & 0.20 & 0.15 & 0.33 \\
11:30 & 0.40 & 0.18 & 0.32 & 0.20 & 0.15 & 0.33 \\
11:45 & 0.39 & 0.19 & 0.32 & 0.20 & 0.15 & 0.33 \\
12:00 & 0.40 & 0.20 & 0.45 & 0.25 & 0.15 & 0.15 \\
12:15 & 0.39 & 0.21 & 0.45 & 0.25 & 0.15 & 0.15 \\
12:30 & 0.39 & 0.22 & 0.45 & 0.25 & 0.15 & 0.15 \\
12:45 & 0.38 & 0.23 & 0.45 & 0.25 & 0.15 & 0.15 \\
13:00 & 0.36 & 0.24 & 0.50 & 0.20 & 0.15 & 0.15 \\
13:15 & 0.36 & 0.24 & 0.50 & 0.20 & 0.15 & 0.15 \\
13:30 & 0.35 & 0.25 & 0.50 & 0.20 & 0.15 & 0.15 \\
13:45 & 0.35 & 0.26 & 0.50 & 0.20 & 0.15 & 0.15 \\
14:00 & 0.35 & 0.27 & 0.40 & 0.20 & 0.25 & 0.15 \\
14:15 & 0.34 & 0.27 & 0.40 & 0.20 & 0.25 & 0.15 \\
14:30 & 0.33 & 0.27 & 0.40 & 0.20 & 0.25 & 0.15 \\
14:45 & 0.32 & 0.27 & 0.40 & 0.20 & 0.25 & 0.15 \\
15:00 & 0.33 & 0.27 & 0.45 & 0.25 & 0.15 & 0.15 \\
15:15 & 0.34 & 0.25 & 0.45 & 0.25 & 0.15 & 0.15 \\
15:30 & 0.35 & 0.23 & 0.45 & 0.25 & 0.15 & 0.15 \\
15:45 & 0.37 & 0.21 & 0.45 & 0.25 & 0.15 & 0.15 \\
16:00 & 0.40 & 0.20 & 0.50 & 0.15 & 0.20 & 0.15 \\
16:15 & 0.37 & 0.25 & 0.50 & 0.15 & 0.20 & 0.15 \\
16:30 & 0.33 & 0.30 & 0.50 & 0.15 & 0.20 & 0.15 \\
16:45 & 0.30 & 0.35 & 0.50 & 0.15 & 0.20 & 0.15 \\
17:00 & 0.27 & 0.40 & 0.50 & 0.20 & 0.15 & 0.15 \\
17:15 & 0.27 & 0.41 & 0.50 & 0.20 & 0.15 & 0.15 \\
17:30 & 0.27 & 0.42 & 0.50 & 0.20 & 0.15 & 0.15 \\
17:45 & 0.26 & 0.43 & 0.50 & 0.20 & 0.15 & 0.15 \\
18:00 & 0.26 & 0.45 & 0.10 & 0.20 & 0.30 & 0.40 \\
18:15 & 0.26 & 0.45 & 0.10 & 0.20 & 0.30 & 0.40 \\
18:30 & 0.26 & 0.45 & 0.10 & 0.20 & 0.30 & 0.40 \\
18:45 & 0.26 & 0.45 & 0.10 & 0.20 & 0.30 & 0.40 \\
19:00 & 0.26 & 0.45 & 0.15 & 0.20 & 0.25 & 0.40 \\
19:15 & 0.26 & 0.45 & 0.15 & 0.20 & 0.25 & 0.40 \\
19:30 & 0.25 & 0.46 & 0.15 & 0.20 & 0.25 & 0.40 \\
19:45 & 0.25 & 0.46 & 0.15 & 0.20 & 0.25 & 0.40 \\
20:00 & 0.25 & 0.47 & 0.18 & 0.12 & 0.25 & 0.45 \\
20:15 & 0.26 & 0.46 & 0.18 & 0.12 & 0.25 & 0.45 \\
20:30 & 0.27 & 0.45 & 0.18 & 0.12 & 0.25 & 0.45 \\
20:45 & 0.28 & 0.44 & 0.18 & 0.12 & 0.25 & 0.45 \\
21:00 & 0.29 & 0.41 & 0.22 & 0.18 & 0.25 & 0.35 \\
21:15 & 0.29 & 0.41 & 0.22 & 0.18 & 0.25 & 0.35 \\
21:30 & 0.29 & 0.41 & 0.22 & 0.18 & 0.25 & 0.35 \\
21:45 & 0.30 & 0.41 & 0.22 & 0.18 & 0.25 & 0.35 \\
22:00 & 0.30 & 0.40 & 0.20 & 0.16 & 0.24 & 0.40 \\
22:15 & 0.27 & 0.40 & 0.20 & 0.16 & 0.24 & 0.40 \\
22:30 & 0.25 & 0.41 & 0.20 & 0.16 & 0.24 & 0.40 \\
22:45 & 0.22 & 0.41 & 0.20 & 0.16 & 0.24 & 0.40 \\
23:00 & 0.21 & 0.42 & 0.15 & 0.20 & 0.20 & 0.45 \\
23:15 & 0.21 & 0.42 & 0.15 & 0.20 & 0.20 & 0.45 \\
23:30 & 0.21 & 0.42 & 0.15 & 0.20 & 0.20 & 0.45 \\
23:45 & 0.21 & 0.42 & 0.15 & 0.20 & 0.20 & 0.45 \\
\bottomrule\end{longtable}


\subsection{Prices, calendar and settlement}
The source is LCCC's IMRP actuals dataset~\cite{lccc_imrp_actuals}, whose reference product is the Great Britain day-ahead hourly price. Source fields identify local date, chronological hourly settlement period and IMRP amount in GBP/MWh. The 2025 series covers 8,760 physical hours. Hourly prices are replicated to four consecutive quarter-hours without smoothing, giving 35,040 committed intervals. The timeline uses UTC for physical ordering and Europe/London for local-date and workload-slot labels. Spring and autumn clock changes produce 92 and 100 quarter-hours in those local days; workload values for skipped/repeated local-clock periods are correspondingly skipped/repeated. No physical UTC interval is missing or duplicated.

Negative prices remain signed in both optimisation and reported costs. There are 672 negative-price quarter-hours in the retained annual data. IMRP is a reference-price signal, not a full retail tariff. The first three hourly prices on 1 January 2026 are GBP~46.49, 40.19 and 28.99/MWh; they supply the twelve final look-ahead intervals. The exact expanded central price timeline is included as \path{data/price_timeline.csv}. Annual cost uses only committed intervals in calendar 2025.
